\section{Introduction}\label{sec:introduction}

Nonconstant adjacency eigenvalues that are small in absolute value make the edge distribution of a
regular graph close to the uniform-density prediction, and thus connect
spectral estimates to expansion \cite{Alon1986}.
The Ramanujan threshold is the sharp asymptotic spectral target in fixed
degree. For a simple $d$-regular graph $G$, its adjacency
matrix $A_G$ has the constant eigenvector $\mathbf1$ with eigenvalue $d$.
The restriction $A_G|_{\mathbf1^\perp}$ measures the remaining spectrum.
Throughout this paper the desired bound is
\[
 \bigl\|A_G|_{\mathbf1^\perp}\bigr\|\le 2\sqrt{d-1}.
\]
Here $2\sqrt{d-1}$ is the spectral radius of the infinite $d$-regular
tree; the Alon--Boppana bound makes it asymptotically optimal
\cite{Alon1986,Nilli1991}. Since $d>2\sqrt{d-1}$ for $d\ge3$, this bound
also forces $G$ to be connected and nonbipartite.

The problem addressed here is to construct such graphs deterministically
at prescribed sizes, without a positive error in the spectral bound.
Our meaning of explicitness is the production of the entire adjacency
list. The degree is fixed throughout the complexity analysis.

\begin{theorem}\label{thm:main}
For every integer $d\ge3$ there are constants $n_0(d)$ and $k_d$ and a
deterministic algorithm with the following property. On input an even
integer $n\ge n_0(d)$, the algorithm outputs the adjacency list of a simple
$d$-regular graph $G$ on $n$ vertices in $O_d(n^{k_d})$ bit operations, and
\[
 -2\sqrt{d-1}<\lambda<2\sqrt{d-1}
 \quad\text{for every eigenvalue $\lambda$ of }A_G|_{\mathbf1^\perp}.
\]
In particular, $G$ is connected and nonbipartite.
\end{theorem}

The exponent and the size threshold may depend on $d$. The statement
concerns all sufficiently large even sizes; it does not assert a
running-time bound polynomial jointly in $n$ and $d$, or a local
adjacency-query algorithm.

\begin{corollary}\label{cor:target-size}
For every fixed $d\ge3$, a deterministic polynomial-time algorithm,
on every sufficiently large integer input $N$, produces a simple
connected nonbipartite $d$-regular Ramanujan graph $G$ with
$N\le |V(G)|\le N+1\le2N$.
\end{corollary}
\begin{proof}
Apply Theorem~\ref{thm:main} to the least even integer $n\ge N$.
\end{proof}

\subsection{Earlier constructions and the exact threshold}

Lubotzky, Phillips and Sarnak and Margulis established arithmetic
constructions attaining the tree threshold in restricted degrees
\cite{LPS1988,Margulis1988}. These constructions obtain their spectral
bounds from arithmetic structure and produce infinite families on
arithmetic sequences of orders. Morgenstern extended the available degrees to $q+1$
for every prime power $q$, including nonbipartite families
\cite{Morgenstern1994}. These exact constructions do not cover every
fixed degree and every prescribed large order simultaneously.

Bilu and Linial introduced a different approach based on graph
two-lifts and signings. A two-lift replaces each vertex by two vertices
and each edge by one of the two matchings between the corresponding
pairs. Its new eigenvalues are those of a signed adjacency matrix,
obtained by assigning a sign $+1$ or $-1$ to each edge.
Their deterministic construction gives an
$O(\sqrt{d\log^3 d})$ bound along iterated lifts
\cite{BL2006}. Marcus, Spielman and Srivastava then used interlacing
polynomials to attain the exact bound for the largest signed eigenvalue.
The bipartite Ramanujan convention excludes both $d$ and $-d$ from the
spectral bound. For bipartite graphs, spectral symmetry also controls the negative edge,
yielding simple bipartite Ramanujan families in every degree
\cite[Theorems~5.3 and~5.5]{MSS2015}.
Hall, Puder and Sawin extended the covering theorem to arbitrary lift
degrees; in particular, simple bipartite $d$-regular Ramanujan graphs
exist on $2dr$ vertices for every positive integer $r$
\cite[Theorem~1.2 and Corollary~2.2]{HPS2018}.

The matching construction of Marcus, Spielman and Srivastava gives
bipartite Ramanujan multigraphs at every even order, with strict bounds
on the eigenvalues other than $\pm d$
\cite[Theorem~1.1]{MSS2018}. Cohen made this construction deterministic
in polynomial time \cite{Cohen2016}. More recently, Cohen and Maor
obtained a quantitative improvement below the tree threshold in this
bipartite multigraph model \cite[Theorem~1.2]{CM2026}.
The restriction to $\mathbf1^\perp$ in Theorem~\ref{thm:main} includes
the negative eigenvalue that bipartiteness would produce.

Recent polynomial expected-time signing algorithms of Jadbabaie,
Saberi and Sra attain the sharp one-sided threshold, and hence the
absolute threshold for bipartite inputs, using a Las Vegas algorithm;
their general two-sided bound has an additional factor $\sqrt2$
\cite[Theorem~6.1]{JSS2026}.
The Bilu--Linial conjecture asks whether every $d$-regular graph has a
two-lift whose new eigenvalues all lie in
$[-2\sqrt{d-1},2\sqrt{d-1}]$ \cite[Conjecture~3.1]{BL2006}.
Our prescribed-order construction uses partial pairings and edge
switches, and does not resolve that signing problem.

For every fixed degree and every fixed $\varepsilon>0$, Mohanty,
O'Donnell and Paredes gave deterministic simple constructions with
bound $2\sqrt{d-1}+\varepsilon$, and order within a factor $1+o(1)$
of a requested size \cite{MOP2020}. Alon obtained this
bound at every sufficiently large size $n$ satisfying $dn$ even
\cite[Theorem~1.3]{Alon2021}. His adjustment of the order deletes
suitably separated vertices from a slightly larger graph and matches
their neighbors; local eigenvector mass estimates control the resulting
spectral error \cite[Section~3]{Alon2021}.
These constructions leave a positive spectral error.
Removing that error requires controlling the eigenvalues on the
much finer scale of their distances from both tree edges.

For every fixed degree and positive error, a uniformly random simple
regular graph is near-Ramanujan with probability tending to one.
This is Friedman's theorem \cite{Friedman2008}; Bordenave gave a new
proof using nonbacktracking operators and extended the method to random
lifts \cite{Bordenave2020}.
Exact existence follows from the joint edge
universality for uniformly random simple regular graphs in every
fixed degree established by Huang, McKenzie and Yau
\cite[Remarks~3.10 and~3.12]{HMY2024}. The probability of the exact two-sided
Ramanujan event tends to approximately $0.69$
\cite[Corollary~1.3]{HMY2024}. Continuity of the
limiting joint edge law gives the same limit for the strict interior
event, so strict existence holds at every sufficiently large
size with $dn$ even. Theorem~\ref{thm:main} supplies a deterministic
polynomial-time construction at the exact threshold and prescribed
even order.

\subsection{How the construction reaches both spectral edges}

Write $q=d-1$, $s=2\sqrt q$ and $e=n^{-2/3+a}$, where $a=10^{-4}$.
The first objective is a simple $d$-regular graph with nonconstant
spectrum in $(-s-e/20,s+e/20)$ and enough local information to repair
the remaining error. An edge switch deletes $\{u,v\}$ and $\{x,y\}$
and inserts $\{u,y\}$ and $\{x,v\}$. It is valid when the four
endpoints are distinct and the inserted edges were absent.
A collection of vertex-disjoint valid switches preserves simplicity,
regularity and the vertex set. A norm bound alone does not determine how
a switch affects an eigenvector concentrated near one of its endpoints.
The construction therefore retains two kinds of local information.

To describe them, let $v_2,\ldots,v_n$ be an orthonormal eigenbasis on
$\mathbf1^\perp$ for this intermediate graph, with eigenvalues
$\lambda_i$, and put $x_i^\sigma=s-\sigma\lambda_i$ for
$\sigma\in\{+1,-1\}$. These signed distances may be slightly negative.
First, the spectral mass near either edge satisfies
\[
 \sum_{i:x_i^\sigma\le w}|v_i(p)|^2\le Cw^{3/2}
 \qquad\text{for every vertex $p$, both signs, and $2e\le w\le2d$}.
\]
Thus the modes requiring repair have little mass at any one vertex.
Summing over vertices also shows that only $O(ne^{3/2})=O(n^{3a/2})$
modes lie within $O(e)$ of either edge. Thus only a small spectral
subspace needs the additional positive bias.
Second, choose one spectrum-avoiding bulk cutoff $w_B^\sigma$ for each sign,
at a prescribed scale larger than $e$, and retain the reciprocal only beyond that cutoff:
\[
 R_B^\sigma=\sum_{i:x_i^\sigma\ge w_B^\sigma}
                  \frac{v_iv_i^T}{x_i^\sigma}.
\]
For all but a small fraction of edges $(p,p')$, the $2$-by-$2$ submatrix
on their endpoints is close, entry by entry, to the tree block
\[
 \frac1{d-2}
 \begin{pmatrix}\sqrt{d-1}&\sigma\\ \sigma&\sqrt{d-1}\end{pmatrix}.
\]
The exceptional edge set may depend on the chosen pair of cutoffs.
This is a local block comparison; the full inverse at the tree edge
need not exist. Theorem~\ref{thm:exact-repair} states the precise scales
and all three repair hypotheses.

The bulk block comparison lets us eliminate the spectral modes beyond
the cutoff and express positivity as a quadratic-form condition on the
remaining near-edge modes. This elimination is a Schur complement.
The original near-edge form can be indefinite. At both edges, a small bias
in the edge-pair distribution gives a positive mean shift to the modes
closest to the threshold; the other near-edge modes retain a fixed
fraction of their positive slack after the mean correction. The local
mass bound controls fluctuations around this mean.
The proof shows that a collection of valid simultaneous switches
satisfying both strict inequalities exists, before giving a finite
deterministic search for it.

\subsection{The four construction stages}

Give each vertex $d$ distinguishable half-edges, called stubs. Pairing
two stubs creates an edge. During the construction we maintain a
positive precision matrix on $\mathbf1^\perp$ for each sign and several
spectral parameters;
at completion these matrices become $zI-\sigma A_G$, with $z>s$ and
$\sigma=\pm1$. Their normalized inverses are compared entry by entry
with truncated weighted sums of nonbacktracking paths, which do not
immediately reverse an edge. Lemma~\ref{lem:discounted-accounts}
selects each pair by minimizing a weighted sum of nonnegative statistics.
Each statistic is discounted by its own computed one-step growth factor,
so the same choice preserves all the required estimates.

\begin{enumerate}
\item \textbf{Pair most stubs.}
Starting from the empty pairing,
Proposition~\ref{prop:early-phase} in Section~\ref{sec:early-phase}
leaves a sparse set of unresolved stubs while preserving positivity,
trace estimates and entrywise comparison. Its main difficulty is simultaneous control
of polynomially many entries over a long sequence of choices.
Lemma~\ref{lem:early-high-moments} bounds the growth of a high moment
with a coefficient independent of the moment order. Taking a
sufficiently large fixed moment therefore gives the required maximum
bound.

\item \textbf{Separate the remaining stubs.}
Given these estimates, Section~\ref{sec:cleanup} pairs the stubs that are close to one another
or to a short return. The surviving stubs occupy distinct vertices,
and no short nonempty nonbacktracking path joins any two of these vertices or
returns to one. The analytic estimates persist during this prescribed
cleanup, by Proposition~\ref{prop:cleanup} and the structural
separation statement in Lemma~\ref{lemma:structure}.

\item \textbf{Complete the graph.}
The separated boundary makes the compression of a smoothly truncated path
reference to unresolved-stub contrasts close to the identity. These contrasts
are vectors of unresolved-stub coordinates whose entries sum to zero. Each vertex row of the
path reference has at most one substantial boundary
coordinate. Exact low-rank inverse updates, given by Woodbury identities,
continue this reference throughout completion.
Section~\ref{sec:boundary-completion} first reduces the squared row
energies and then uses those smaller energies to improve the entry
bounds. This order permits an arbitrary matching of the last few stubs.
Proposition~\ref{prop:boundary-completion} supplies a simple regular
graph with the required spectral confinement and entry comparison
against one smoothly parameterized reference. The construction retains
the path information at the separation stop needed to analyze that
same reference.

\item \textbf{Remove the spectral excess.}
In Section~\ref{sec:exact-repair}, the parameter dependence of the frozen
reference gives the local mass bound in Lemma~\ref{lem:repair-projector};
the retained path information gives the typical endpoint blocks in
Lemma~\ref{lem:repair-bulk}. Together with confinement, these are the
three inputs to Theorem~\ref{thm:exact-repair}. Its biased switches
make both signed precisions strictly positive at $s$ on
$\mathbf1^\perp$, while preserving simplicity, regularity and the
vertex set.
\end{enumerate}

Figure~\ref{fig:construction} summarizes these successive outputs.
\begin{figure}[hbp]
\centering
\begin{tikzpicture}[
  box/.style={draw=MidnightBlue!65,fill=MidnightBlue!3,rounded corners=2pt,
    text width=.39\textwidth,minimum height=23mm,inner sep=8pt,align=left,font=\small},
  flow/.style={-{Latex[length=2mm]},thick,draw=MidnightBlue!80}]
\node[box] (early) {\textbf{1. Early pairing}\\
  Positive signed precisions;\\trace and entry bounds.};
\node[box,right=8mm of early] (clean) {\textbf{2. Separation}\\
  Distinct unresolved vertices;\\no short joins or returns.};
\node[box,below=10mm of clean] (complete) {\textbf{3. Boundary completion}\\
  Small row energies, then entry errors;\\nonconstant spectrum in\\
  \mbox{$(-s-e/20,s+e/20)$}.};
\node[box,left=8mm of complete] (repair) {\textbf{4. Local estimates and switches}\\
  Near-edge mass and typical bulk blocks;\\nonconstant spectrum in $(-s,s)$.};
\draw[flow] (early)--(clean);
\draw[flow] (clean)--(complete);
\draw[flow] (complete)--(repair);
\end{tikzpicture}
\caption{The four stages on the same $n$ vertices. The first three
produce the graph and a reference matrix from which the local spectral
estimates for repair are extracted.
Paths in the separation stage are nonempty nonbacktracking paths.
All precisions act on $\mathbf1^\perp$; positivity is tracked for both
spectral signs throughout.}
\label{fig:construction}
\end{figure}


The deterministic selection rule belongs to the method of conditional
expectations and pessimistic estimators \cite{Raghavan1988,WX2008}.
Here each statistic has its own discount, and the transition estimates
control those discounts along the selected history. Weighting a path of
length $j$ by $(\sigma h)^j$ introduces a scalar variable $h$. For a
completed graph, the nonbacktracking path recurrence leads to the quadratic matrix
$I-\sigma hA_G+(d-1)h^2I$, the regular-graph Ihara--Bass pencil
\cite[Sections~6 and~8]{FZ1999}; the partial configuration instead
averages over unresolved stubs, so its identities and quantitative
estimates are proved directly. Likewise, degree-preserving switchings
are classical in regular-graph generation \cite{MW1990}; the local
spectral inputs and the bias that improves both edges are established
in the repair proof.

Local comparison with tree Green functions and stability under edge
resampling were developed by Bauerschmidt, Huang and Yau for random
regular graphs of sufficiently large fixed degree in the bulk
\cite[Sections~2, 7 and~11]{BHY2019}. Huang and Yau extended this
local-resampling approach to every fixed degree $d\ge3$ and obtained
control up to the spectral edges \cite[Sections~1.4, 3 and~5]{HY2024}.
These are methodological precedents for the local comparisons used
here; the estimates along the deterministic pairing history and the
biased simultaneous repair are proved below.

Section~\ref{sec:setup} fixes notation and parameter order;
Section~\ref{sec:accounts} gives the selection principle; and
Section~\ref{sec:foundations} develops the configuration identities
and structural estimates. Section~\ref{sec:algorithm} assembles the
four stages and proves polynomial bit complexity. Its final check is
exact rational positivity of $B^T(4(d-1)I-A_G^2)B$, where the columns
of $B$ form a rational basis of $\mathbf1^\perp$; the repair theorem
ensures that the polynomial search contains a passing candidate.
The structural
enumeration uses color-coding \cite{AYZ1995} and the perfect-hash
families of Naor, Schulman and Srinivasan
\cite[Theorem~3(iii)]{NSS1995}; the weighted configuration estimates
and spectral repair are proved here.
